\chapter{Observabilidad y seguridad}
\label{chap:observabilidad-seguridad}

\section{Comprobaciones de salud}
\label{sec:health-checks}

La API expone dos endpoints de salud. \path{/api/v1/health/live} indica que el proceso está vivo. \path{/api/v1/health/ready} comprueba la conectividad real con PostgreSQL mediante el \texttt{DbContext}. Esta distinción permite diferenciar entre una aplicación arrancada y una aplicación preparada para servir tráfico.

En AWS, el target group del Application Load Balancer utiliza la ruta de readiness para decidir si la tarea ECS puede recibir tráfico. También se ha usado el mismo endpoint como prueba de humo final. En Sprint 6 se fijó esta ruta como valor por defecto parametrizable del target group y se añadió un periodo de gracia de ECS para evitar que tareas recién arrancadas sean sustituidas antes de completar el arranque.

\section{Logs y métricas}
\label{sec:logs-metricas}

Los logs del contenedor se envían a CloudWatch Logs mediante el driver \texttt{awslogs} de ECS. El grupo de logs sigue la convención \path{/aws/ecs/transitops/dev/api}. Esta integración permite inspeccionar errores de arranque, fallos de configuración, problemas de conexión con RDS y salidas de tareas puntuales como las migraciones.

En Sprint 6 se sustituyó la salida de consola no estructurada por logging JSON mediante \texttt{AddJsonConsole}. La configuración incluye scopes, marcas temporales UTC y campos que CloudWatch puede filtrar. También se añadió un middleware de correlación de peticiones:

\begin{itemize}
  \item acepta \texttt{X-Correlation-ID} si el cliente lo proporciona;
  \item genera un identificador nuevo si el header no existe;
  \item devuelve siempre \texttt{X-Correlation-ID} en la respuesta;
  \item alinea \texttt{HttpContext.TraceIdentifier} con ese valor;
  \item abre un scope de logging con \texttt{CorrelationId};
  \item registra por petición método, ruta, código HTTP, duración, usuario y rol cuando existen.
\end{itemize}

Esto permite seguir una petición concreta desde la respuesta HTTP hasta los logs de CloudWatch. Además, los errores gestionados por el middleware de excepciones heredan el mismo identificador, por lo que el campo de error devuelto al cliente y el log del servidor quedan correlados.

\begin{figure}[htbp]
\centering
\fbox{\begin{minipage}{0.88\textwidth}
\centering
\vspace{1.2cm}
\textbf{Placeholder de captura}\\
Consulta de CloudWatch Logs Insights filtrando por \texttt{State.CorrelationId = "sprint6-login-final"} y mostrando logs JSON de una petición de login.
\vspace{1.2cm}
\end{minipage}}
\caption{Evidencia prevista de logs JSON correlables.}
\label{fig:placeholder-cloudwatch-logs-correlation}
\end{figure}

Las métricas básicas de ECS, ALB y RDS quedan disponibles en CloudWatch como parte de los servicios gestionados. El módulo de observabilidad añade un dashboard \path{transitops-dev-api} con paneles para ALB, utilización de ECS, RDS y logs recientes. El dashboard no sustituye a una plataforma APM completa, pero ofrece una superficie operativa suficiente para un entorno \texttt{dev} defendible.

\section{Alarmas y operación}
\label{sec:alarmas-operacion}

La operación de Sprint 6 deja de depender solo de comprobaciones manuales. Terraform crea alarmas CloudWatch para:

\begin{itemize}
  \item errores de aplicación detectados por filtro métrico sobre logs JSON;
  \item respuestas 5xx del target group del ALB;
  \item tiempo de respuesta del ALB;
  \item ausencia de targets sanos;
  \item utilización de CPU de ECS;
  \item utilización de memoria de ECS;
  \item CPU de RDS;
  \item conexiones de base de datos;
  \item almacenamiento libre en RDS.
\end{itemize}

Las alarmas usan \texttt{treat\_missing\_data} conservador para reducir ruido cuando \texttt{dev} está destruido o con \texttt{desired\_count = 0}. Si se configura \texttt{alarm\_email}, Terraform crea un topic SNS y una suscripción de email; si no se configura, las alarmas existen sin acciones. Durante la validación de Sprint 6 no se creó topic SNS porque no se proporcionó email.

\begin{figure}[htbp]
\centering
\fbox{\begin{minipage}{0.88\textwidth}
\centering
\vspace{1.2cm}
\textbf{Placeholder de captura}\\
Dashboard CloudWatch \path{transitops-dev-api} con widgets de ALB, ECS, RDS y logs recientes.
\vspace{1.2cm}
\end{minipage}}
\caption{Evidencia prevista del dashboard operativo.}
\label{fig:placeholder-cloudwatch-dashboard}
\end{figure}

\section{Seguridad aplicada}
\label{sec:seguridad-aplicada}

Las medidas de seguridad aplicadas son:

\begin{itemize}
  \item Autenticación mediante login y emisión de \acrshort{jwt}.
  \item Contraseñas almacenadas como hash, no en claro.
  \item Autorización por roles \texttt{admin} y \texttt{operator}.
  \item Bootstrap del primer administrador protegido por token externo.
  \item Secretos en Secrets Manager, no en el repositorio.
  \item RDS en subredes privadas y sin acceso público.
  \item ECS en subredes privadas, recibiendo tráfico solo desde ALB.
  \item ALB como único punto de entrada público.
  \item HTTPS mediante ACM y redirección HTTP a HTTPS.
  \item GitHub Actions con OIDC en lugar de claves AWS permanentes.
  \item Correlación de errores y respuestas mediante \texttt{X-Correlation-ID}, evitando exponer trazas internas al cliente.
\end{itemize}

\section{Limitaciones de seguridad}
\label{sec:limitaciones-seguridad}

El sistema no pretende cubrir todos los controles exigibles a un producto crítico en producción. Quedan fuera del alcance actual la rotación automática de secretos, WAF, protección avanzada contra abuso, auditoría detallada de acciones administrativas, cifrado con claves KMS propias, políticas IAM de mínimo privilegio refinadas al extremo, pruebas de penetración, OpenTelemetry/X-Ray y hardening exhaustivo de cabeceras HTTP.

Estas limitaciones no invalidan el objetivo del trabajo, ya que se ha priorizado una línea base realista y demostrable: autenticación, autorización, red privada, secretos externos, HTTPS e infraestructura reproducible.

\section{Revisión de seguridad Sprint 8}
\label{sec:revision-seguridad-sprint8}

Sprint 8 añadió una revisión explícita de seguridad sobre la arquitectura desplegable. La conclusión principal es que la aplicación mantiene una frontera pública pequeña: el ALB es el único punto de entrada, HTTP se usa únicamente para redirección a HTTPS, ECS no recibe IP pública y RDS no es accesible desde Internet.

En IAM se diferencian tres casos. El execution role de ECS puede leer únicamente los secretos referenciados por la task definition. El task role de la aplicación no tiene permisos AWS porque la API no llama a servicios AWS en ejecución. El rol OIDC de GitHub conserva permisos amplios sobre los servicios gestionados por Terraform para poder aplicar y destruir el entorno \texttt{dev}; se acepta como riesgo controlado porque la trust policy queda limitada al repositorio y rama configurados, y porque el entorno se destruye tras capturar evidencia.

\begin{figure}[htbp]
\centering
\fbox{\begin{minipage}{0.88\textwidth}
\centering
\vspace{1.2cm}
\textbf{Placeholder de captura}\\
Ejecución del script \path{Test-Sprint8AwsPosture.ps1} mostrando RDS privado, ECS sin IP pública, grupos de seguridad esperados, secretos en Secrets Manager y advertencia documentada sobre el rol OIDC de Terraform.
\vspace{1.2cm}
\end{minipage}}
\caption{Evidencia de revisión de seguridad Sprint 8.}
\label{fig:placeholder-sprint8-security-review}
\end{figure}
